\section{结论}

本文提出了\textbf{EdgeSched}，一个专为资源受限边缘设备上多模态AI任务共存而设计的轻量级资源调度框架。通过解决"8GB显存墙"问题，EdgeSched使得常驻LLM助手与突发图像生成任务能够并发运行，而无需昂贵的硬件升级。

我们的主要贡献包括：
\begin{enumerate}
\item \textbf{非对称资源迁移}：提出了一种结合"紧急模式"(Urgent Mode)用于快速抢占(GPU→CPU)与"原子热切换"(Atomic Hot-Swap)用于无缝恢复(CPU→GPU)的新颖策略，在显存安全与服务连续性之间取得了平衡。

\item \textbf{轻量级持久化}：基于mmap的软件定义状态背板，提供了可忽略开销(< 0.01ms)的崩溃一致性，消除了系统故障后的"僵尸任务"。

\item \textbf{可预测QoS}：结合优先级队列、背压机制和取消语义，确保系统即使在饱和状态下仍能保持响应和稳定。
\end{enumerate}

在消费级硬件(RTX 5070 Laptop)上的实验结果验证了我们的方法。在基线方法因OOM而失败的压力测试中，EdgeSched实现了\textbf{100\%的成功率}，在模型恢复期间保持了\textbf{低于300ms的服务中断}，并在\textbf{15ms}内从崩溃中恢复。虽然原始吞吐量略低于简单的信号量方案(~11\%)，但在可靠性和用户体验方面的收益使得EdgeSched成为下一代边缘AI操作系统的可行基础。

\subsection{未来工作}

未来的工作将集中在以下几个方向：

\begin{enumerate}
\item \textbf{异构多GPU支持}：扩展EdgeSched以支持多GPU配置，实现更细粒度的资源分配和负载均衡。

\item \textbf{内核级调度}：集成细粒度的内核级调度以实现更低延迟的抢占，进一步减少服务中断时间。

\item \textbf{自适应策略}：基于历史负载模式和用户行为，动态调整迁移策略和队列容量，实现更智能的资源管理。

\item \textbf{更多模态支持}：扩展系统以支持语音合成(TTS)、视觉-语言模型(VL)等更多模态任务，构建完整的多模态AI助手生态。
\end{enumerate}

EdgeSched为边缘设备上的多模态AI应用提供了一个实用的调度解决方案，我们期待它能够推动边缘智能领域的进一步发展。
